\section{Contexto del Proyecto y Objetivos de la Practica}

\subsection{Contexto tecnico de partida}
El proyecto \textit{Didascalia VC} se encontraba en fase de evolucion funcional, con una base desarrollada en Unity y dependencias de XR para ejecucion en entornos de Realidad Virtual. El producto debia consolidar estabilidad tecnica y escalabilidad de comportamientos antes de su ampliacion con nuevos casos de uso pedagogicos.

Al inicio de las practicas se identificaron, entre otras, las siguientes necesidades:
\begin{itemize}[leftmargin=1.5em]
    \item Mejorar la fluidez en VR, reduciendo picos de carga y tirones perceptibles.
    \item Definir un marco extensible para comportamientos complejos de estudiantes.
    \item Endurecer la capa de conexion con la web para detectar y gestionar mensajes no esperados.
    \item Verificar la cobertura completa de conflictos soportados y ampliar el catalogo con nuevos escenarios.
    \item Garantizar despliegue estable en Meta Quest 3.
\end{itemize}

\subsection{Objetivos de aprendizaje y contribucion}
La asignacion de practicas estuvo alineada con un perfil tecnico orientado a desarrollo de software interactivo. Se definieron como objetivos personales y de contribucion:
\begin{enumerate}[leftmargin=1.5em]
    \item Aplicar metodologias de diagnostico y optimizacion en un entorno VR real.
    \item Dise\~nar estructuras de datos y flujo de estados mantenibles para crecimiento funcional.
    \item Implementar validaciones y aserciones para reforzar fiabilidad de ejecucion.
    \item Integrar cambios de codigo, contenido y configuracion manteniendo compatibilidad entre editor, build y dispositivo final.
\end{enumerate}

\subsection{Tecnologias y herramientas empleadas}
\begin{itemize}[leftmargin=1.5em]
    \item Motor: Unity (proyecto C\# con XR Interaction Toolkit y configuracion para Oculus/Meta).
    \item Lenguaje principal: C\#.
    \item Control de versiones: Git (rama de trabajo \textit{Didascalia-Remaster}).
    \item Entorno de pruebas: escenas de menu y aula, validacion de eventos y estados en tiempo de ejecucion.
    \item Objetivo hardware: Meta Quest 3 y entorno de desarrollo de escritorio para depuracion.
\end{itemize}

\subsection{Metodologia de trabajo}
La metodologia seguida ha sido incremental, en ciclos cortos de implementacion y verificacion. Cada bloque tecnico se abordo con una secuencia de: analisis del problema, desarrollo de solucion, pruebas en escena/build y ajustes de integracion para minimizar regresiones.

Este enfoque se refleja en el historial de commits del periodo analizado (ultimos cinco meses), donde se observan iteraciones continuas sobre rendimiento, sistema de estados/comportamientos, red/web y adaptacion de contenido.